\chapter{Validation}\label{ch:ml:validation}

Le modèle d'une équipe de recherche affiche un \omterm{def:ml:why-financial-machine-learning-is-different:r2}{R² hors échantillon} de 1.3~\% sur les sept dernières années d'un panel mensuel d'actions,
soit trois fois ce qu'atteignent les meilleurs modèles publiés. L'analyse a posteriori ne relève aucune erreur dans le modèle~: la
surprise de résultats qu'il utilise a été rattachée à chaque action dès la fin du trimestre comptable, et le rapport du trimestre parvient au marché
un mois plus tard, durant le mois même dont le modèle prédit le rendement. Sur les mêmes données, rattachées à la date de dépôt, le modèle
obtient $-0.19~\%$, comme il se doit~: dans le monde de ce chapitre, rien n'est prévisible au moment de la décision, si bien que tout
nombre positif ci-dessous est une fuite. La validation estime comment un modèle se comportera sur des données qu'il n'a pas vues~; ce chapitre
traite des nombreuses façons dont les données qu'il n'a pas vues finissent malgré tout par y entrer, et de la manière de repérer chacune d'elles.

\section{Ce qu'estime la validation}

La validation répond à deux questions différentes, et un seul nombre ne peut répondre aux deux.

\begin{definition}[Sélection de modèle, recherche d'hyperparamètres]\label{def:ml:validation:selection}
La \emph{sélection de modèle}\index{sélection de modèle} choisit, parmi les modèles ou configurations candidats, celui qui est censé
généraliser le mieux. Une \emph{recherche d'hyperparamètres}\index{recherche d'hyperparamètres} est une sélection de modèle sur un ensemble de
valeurs d'\omterm{def:ml:why-financial-machine-learning-is-different:sets}{hyperparamètres}~: une grille, des tirages aléatoires (Bergstra et Bengio, 2012), ou une recherche séquentielle qui propose de nouvelles
valeurs à partir des scores des précédentes.
\end{definition}

La seconde question est celle de l'\emph{évaluation}~: quelle sera la performance du modèle retenu~? Le score qui a remporté la sélection est une
réponse biaisée à cette question.

\begin{proposition}[Le score du vainqueur est optimiste]\label{prop:ml:validation:max}
Soient $\hat S_1,\dots,\hat S_K$ des estimateurs sans biais des scores de généralisation $S_1,\dots,S_K$ de $K$ candidats. Alors
$\E[\max_k\hat S_k] \ge \max_kS_k$, avec égalité seulement si le maximum est atteint par le même candidat dans presque chaque
tirage, et sans erreur d'estimation.
\end{proposition}

\begin{proof}
$\max$ est convexe, donc par l'inégalité de Jensen $\E[\max_k\hat S_k] \ge \max_k\E[\hat S_k] = \max_kS_k$.
\end{proof}

Le biais croît avec le nombre de candidats et avec le bruit de chaque score, lequel est important sur les données de marché. Le livre~4
(chapitre~12) a transformé cela en ratio de Sharpe déflaté~; ici, c'est la raison pour laquelle un score de sélection ne doit jamais être présenté
comme une évaluation.

\section{Découpages qui respectent le temps}

L'univers du chapitre est le panel de \texttt{firm.mlsynth}, dont le \omterm{def:ml:why-financial-machine-learning-is-different:r2}{plafond de prévisibilité} est nul~: 200 actions, 240 mois, vingt
caractéristiques qui ne contiennent aucune information. Chaque action publie ses résultats tous les trimestres~; le rapport du trimestre clos au mois $t$ est déposé
pendant le mois $t+1$ et déplace le rendement de ce mois de 2~\% par unité de surprise standardisée. Le pipeline est rigoureux
partout ailleurs~: des arbres de gradient boosting entraînés sur des rendements centrés chaque mois, sur les 160 premiers mois moins un, et évalués
par rapport à zéro sur les 80 derniers.

Le livre~7 (chapitre~20) a construit les découpages qui respectent le temps~: la validation $k$-fold purgée avec embargo, la validation croisée combinatoire
purgée et l'analyse glissante (walk-forward). \texttt{firm.cvsplit} les enveloppe en validateurs scikit-learn, si bien que la
bibliothèque~: ses outils de recherche et de scoring les utilisent sans modification (\cref{fig:ml:validation:splits}). Le chevauchement des étiquettes
qu'ils préviennent est sévère sur les panels~: avec des étiquettes à trois mois échantillonnées chaque mois et une validation croisée à 4 plis, un pli mélangé laisse
38\,080 étiquettes d'entraînement qui partagent un mois avec une étiquette de test~; le pli purgé n'en laisse aucune. Les arbres profonds sur les vingt
caractéristiques de bruit obtiennent 3.47~\% avec des plis mélangés et $-3.37~\%$ avec des plis purgés. Les caractéristiques sont
persistantes, si bien qu'un modèle souple peut reconnaître une action une année donnée et en déduire les étiquettes des mois voisins.
Un classifieur chargé de distinguer, à partir des seules caractéristiques, la première moitié des mois de la seconde atteint une AUC
de 0.71~: les variables localisent le temps, et c'est là l'avertissement.

\begin{omfigure}
\begin{tikzpicture}[x=0.69cm,y=0.55cm,font=\footnotesize]
\def\row#1#2{\node[anchor=east] at (-0.2,#1) {#2};}
\row{4}{$k$-fold mélangé}
\foreach \i/\c in {0/1,1/0,2/0,3/1,4/0,5/0,6/1,7/0,8/0,9/1,10/0,11/0,12/1,13/0,14/0,15/1} {
  \ifnum\c=1 \fill[omThm!45] (\i,3.7) rectangle ++(0.92,0.6); \else \fill[omDef!30] (\i,3.7) rectangle ++(0.92,0.6); \fi }
\row{3}{$k$-fold purgé, embargo}
\foreach \i in {0,1,2,3,4,10,11,12,13,14,15} { \fill[omDef!30] (\i,2.7) rectangle ++(0.92,0.6); }
\foreach \i in {6,7,8} { \fill[omThm!45] (\i,2.7) rectangle ++(0.92,0.6); }
\fill[gray!25] (5,2.7) rectangle ++(0.92,0.6); \fill[gray!25] (9,2.7) rectangle ++(0.92,0.6);
\row{2}{walk-forward}
\foreach \i in {2,3,4,5,6,7,8,9} { \fill[omDef!30] (\i,1.7) rectangle ++(0.92,0.6); }
\foreach \i in {10,11} { \fill[omThm!45] (\i,1.7) rectangle ++(0.92,0.6); }
\row{1}{imbriquée (un pli externe)}
\foreach \i in {0,1,2,3,4,5,6,7,8,9} { \fill[omDef!30] (\i,0.7) rectangle ++(0.92,0.6); }
\foreach \i in {3,4,5} { \fill[omProp!55] (\i,0.7) rectangle ++(0.92,0.6); }
\foreach \i in {12,13,14,15} { \fill[omThm!45] (\i,0.7) rectangle ++(0.92,0.6); }
\fill[gray!25] (10,0.7) rectangle ++(0.92,0.6); \fill[gray!25] (11,0.7) rectangle ++(0.92,0.6);
\draw[->] (0,0.1) -- (16.3,0.1) node[right]{temps};
\fill[omDef!30] (0.2,-0.9) rectangle ++(0.5,0.4); \node[anchor=west] at (0.8,-0.7) {apprentissage};
\fill[omThm!45] (4.1,-0.9) rectangle ++(0.5,0.4); \node[anchor=west] at (4.7,-0.7) {test};
\fill[omProp!55] (6.1,-0.9) rectangle ++(0.5,0.4); \node[anchor=west] at (6.7,-0.7) {validation interne};
\fill[gray!25] (10.9,-0.9) rectangle ++(0.5,0.4); \node[anchor=west] at (11.5,-0.7) {purgé ou sous embargo};
\end{tikzpicture}

\omcaption{Quatre façons de découper seize périodes. Les plis mélangés dispersent les périodes de test parmi les périodes d'entraînement, si bien que
les étiquettes qui se chevauchent et les variables persistantes fuient~; la purge retire les étiquettes d'entraînement qui chevauchent le bloc de test et
l'embargo juste après~; l'analyse glissante n'entraîne que sur le passé~; la \omterm{def:ml:validation:nested}{validation croisée imbriquée} choisit les \omterm{def:ml:why-financial-machine-learning-is-different:sets}{hyperparamètres} sur
les plis internes des seules données d'entraînement externes.}\label{fig:ml:validation:splits}
\end{omfigure}

\section{Choisir les hyperparamètres sans s'illusionner}

\begin{definition}[Validation croisée imbriquée]\label{def:ml:validation:nested}
La \emph{validation croisée imbriquée}\index{validation croisée imbriquée} exécute une validation croisée interne dans chaque \omterm{def:ml:why-financial-machine-learning-is-different:sets}{ensemble d'entraînement}
d'une validation croisée externe~: les plis internes choisissent les \omterm{def:ml:why-financial-machine-learning-is-different:sets}{hyperparamètres}, la configuration ainsi retenue est réentraînée sur l'\omterm{def:ml:why-financial-machine-learning-is-different:sets}{ensemble
d'entraînement} externe et évaluée sur le pli de test externe. Le score moyen externe évalue l'ensemble de la procédure, sélection comprise.
\end{definition}

Le bruit de sélection se voit le plus facilement là où il n'y a rien à sélectionner. Trente exécutions d'une même configuration d'arbres boostés,
identiques hormis leur graine aléatoire, sont entraînées sur les 159 premiers mois du chapitre. Leur R² sur les mois 160
à 199 s'étend de $-0.60~\%$ à $-0.27~\%$, avec une médiane de $-0.40~\%$~; la meilleure d'entre elles paraît supérieure de 0.13 points à
la médiane. Sur les mois 200 à 239, que le choix n'a jamais vus, cette exécution obtient $-0.20~\%$ et l'exécution médiane $-0.19~\%$
(\cref{fig:ml:validation:seeds})~: sur les trente graines, les scores des deux périodes sont corrélés à $-0.16$. Les 0.13 points
étaient le bruit propre de la sélection, comme l'énonce la \cref{prop:ml:validation:max}.

\begin{omfigure}
\begin{tikzpicture}
\begin{axis}[width=8.6cm,height=6.2cm,xlabel={R² sur les mois de sélection (\%)},ylabel={R² sur les mois intacts (\%)},
  tick label style={font=\small},label style={font=\small},xmin=-0.65,xmax=-0.2,ymin=-0.35,ymax=0.18,grid=major,
  grid style={gray!25},table/col sep=comma,
  xticklabel style={/pgf/number format/fixed,/pgf/number format/precision=1},scaled x ticks=false]
\addplot[only marks,mark=*,mark size=1.6pt,omDef] table[x=selection,y=fresh]{figdata/ml/03-validation/seeds.csv};
\draw[omThm,thick] (axis cs:-0.2723,-0.2) circle[radius=4pt];
\node[font=\footnotesize,omThm,anchor=south] at (axis cs:-0.2723,-0.165) {retenu};
\end{axis}
\end{tikzpicture}

\omcaption{Trente graines des mêmes arbres boostés sur un monde où il n'y a rien à prédire~: le score sur les mois servant à
les départager, comparé au score sur les mois que personne n'a regardés. La graine retenue pour son score de sélection
(entourée) reste ordinaire ensuite. Données~: \texttt{ml\_validation.selection}.}\label{fig:ml:validation:seeds}
\end{omfigure}

La \omterm{def:ml:validation:nested}{validation croisée imbriquée} corrige le chiffre annoncé. Avec huit graines de ce type comme candidats, trois plis externes purgés et trois
internes sur un mois sur deux (pour que les étiquettes ne se chevauchent pas), la recherche à plat annonce son meilleur score externe moyen,
$-0.62~\%$~; la procédure imbriquée, qui choisit à l'intérieur de chaque \omterm{def:ml:why-financial-machine-learning-is-different:sets}{ensemble d'entraînement} externe, obtient $-0.85~\%$. Le chiffre honnête est le
pire. Les candidats étaient des graines, ce qui rend le propos évident~; avec de vrais \omterm{def:ml:why-financial-machine-learning-is-different:sets}{hyperparamètres}, le même biais se superpose à
tout écart que produisent les \omterm{def:ml:why-financial-machine-learning-is-different:sets}{hyperparamètres}.

\begin{method}[Une recherche qui peut être rapportée]\label{met:ml:validation:search}
\begin{enumerate}
\item Fixer les découpages (plis purgés ou walk-forward, à partir des intervalles des étiquettes) et le score avant la première exécution.
\item Chercher uniquement sur les plis internes~; consigner chaque configuration essayée dans le journal de recherche (livre~7, chapitre~1)~: le
nombre d'essais fait partie du résultat.
\item Rapporter le score imbriqué (ou walk-forward) de la procédure, jamais le meilleur score interne obtenu.
\item Conserver un échantillon final réservé que nulle recherche n'a touché, et ne l'utiliser qu'une fois.
\end{enumerate}
\end{method}

\section{Chasse aux fuites}

\begin{definition}[Fuite de la cible, contamination entraînement--test]\label{def:ml:validation:leakage}
La \emph{fuite de la cible}\index{fuite de la cible} (target leakage) est une information contenue dans une variable, ou dans une étape de la chaîne de traitement, qui ne serait pas
disponible à l'instant de décision que la prédiction représente, typiquement parce qu'elle a été dérivée de la cible
ou enregistrée après elle. La \emph{contamination entraînement--test}\index{contamination entraînement--test} (train--test contamination) est une information sur les données de test
qui parvient à l'ajustement par les données elles-mêmes~: étiquettes qui se chevauchent ou enregistrements dupliqués des deux côtés d'une séparation,
ou des statistiques calculées sur l'ensemble de l'échantillon.
\end{definition}

Le livre~7 a nommé le biais d'anticipation (look-ahead bias, chapitre~3) et le biais de survie~; la fuite en est la forme générale, et la littérature sur
la fouille de données y voit l'une des erreurs les plus fréquentes du domaine (Kaufman, Rosset, Perlich et Stitelman). Le chapitre
introduit cinq fuites dans sa chaîne de traitement soignée (\cref{tab:ml:validation:leaks}). Deux sont dans les variables~: la jointure de fin de période
de la surprise, et \emph{encodage par la cible} (target encoding), ici le rendement moyen du secteur de l'action le même mois, une statistique
de la cible qui se trouve être disponible dans la table au moment où le modèle est ajusté. Une se trouve dans la sélection de variables~:
le criblage de 1\,000 candidats de pur bruit sur leur coefficient d'information sur l'échantillon entier avant la séparation, puis en retenant
dix. Une se trouve dans les découpages (plis mélangés sur des étiquettes qui se chevauchent), et une dans le choix (la meilleure graine au score de test).

\begin{table}[htbp]
\centering\small
\begin{tabular}{@{}>{\raggedright\arraybackslash}p{4.2cm}>{\centering\arraybackslash}p{2.3cm}>{\centering\arraybackslash}p{2.3cm}>{\raggedright\arraybackslash}p{5.2cm}@{}}
\toprule
fuite & $R^2_{\mathrm{oos}}$ avec fuite & version honnête & détectée par\\
\midrule
jointure de fin de période de la surprise & 1.27~\% & $-0.19~\%$ (date de dépôt) & test de troncature\\
encodage par la cible, même mois & 8.97~\% & $-0.19~\%$ (aucun) & canari, test de troncature\\
criblage sur tout l'échantillon & 0.13~\% & $-0.16~\%$ (mois d'entraînement) & canari, test de troncature\\
plis mélangés, étiquettes qui se chevauchent & 3.47~\% & $-3.37~\%$ (plis purgés) & contrôle de chevauchement des plis\\
meilleure des 30 graines sur le test & $-0.27~\%$ & $-0.20~\%$ (mois non touchés) & échantillon réservé, recherche imbriquée\\
\bottomrule
\end{tabular}
\caption{Cinq fuites dans un monde où rien n'est prévisible à l'instant de décision (200 actions, 240 mois). Tout chiffre positif
est faux~; la fuite de la dernière ligne représente un gain de 0.13 point par rapport à la graine médiane, gain qui ne se répète pas. Le criblage utilise
une régression ridge sur les dix variables criblées, les autres des arbres boostés. Données~: \texttt{ml\_validation.leaks}.}\label{tab:ml:validation:leaks}
\end{table}

\begin{definition}[Validation adverse, audit de fuite]\label{def:ml:validation:audit}
La \emph{validation adverse}\index{validation adverse} (adversarial validation) entraîne un classifieur à distinguer les lignes d'entraînement des lignes de test~; une
AUC bien supérieure à un demi signifie que les variables permettent de repérer la séparation, donc que les plis mélangés fuiront et que la période de test diffère
de la période d'entraînement. Un \emph{audit de fuite}\index{audit de fuite} (leakage audit) est l'ensemble des contrôles qu'une chaîne de traitement passe avant que son
score soit jugé crédible~: au moins un test de troncature de chaque variable, un passage canari de toute la chaîne de traitement, un contrôle de chevauchement des
découpages, et la validation adverse.
\end{definition}

Le test de troncature recalcule chaque variable sur les données connues à une date donnée et la compare à la valeur enregistrée~: sur la
surprise enregistrée par mois de dépôt, la variable de la jointure de fin de période diffère d'au plus 4.77 écarts-types, celle de la jointure à la date de dépôt
diffère de zéro exactement~; les variables criblées diffèrent de 5.13, car le criblage a retenu d'autres candidats avec moins
de données. Le canari fait tourner toute la chaîne de traitement sur une cible qui n'est que du bruit par construction, avec la forme de la vraie cible~: la
chaîne à encodage par la cible obtient encore 4.03~\% sur du bruit (erreur type 0.02~\%) et la chaîne de criblage 0.17~\% (0.03~\%),
ce qu'aucune chaîne honnête ne peut produire. Le canari ne détecte pas la jointure de fin de période ($-0.44~\%$), car cette variable fuit
sur la vraie cible, et non sur la chaîne elle-même~: l'audit a besoin des deux tests. Le contrôle de chevauchement des découpages compte les étiquettes d'entraînement
dont les intervalles touchent le bloc de test~; l'échantillon réservé détecte la sélection. Aucun contrôle isolé ne détecte les cinq.

\begin{method}[L'audit de fuite]\label{met:ml:validation:audit}
\begin{enumerate}
\item Stocker chaque entrée avec son temps de connaissance (le magasin bitemporel du livre~7) et recalculer chaque variable sur les données
connues à un échantillon de dates~; toute différence signale une fuite.
\item Faire tourner toute la chaîne de traitement sur une cible canari ayant la structure de la vraie cible (même chevauchement, même échelle)~; un score
dépassant le hasard de plus de deux erreurs types signale une fuite dans la chaîne de traitement.
\item Vérifier qu'aucune étiquette d'entraînement ne chevauche une étiquette de test~; lancer la \omterm{def:ml:validation:audit}{validation adverse} sur les variables.
\item Conserver un échantillon réservé en dehors de toute recherche, et compter les essais.
\end{enumerate}
\end{method}

\section{Tutoriel~: cinq fuites}\label{tut:ml:validation:tut}

\begin{tutorial}
\textbf{Objectif.} Introduire cinq fuites dans une chaîne de traitement sur un monde où il n'y a rien à prédire, mesurer le R² que chacune fabrique,
détecter chacune avec un détecteur, et comparer les recherches à plat et imbriquées. \textbf{État final~:} le \cref{tab:ml:validation:leaks},
la \cref{fig:ml:validation:seeds}, les valeurs des détecteurs.
\begin{tutsteps}
\item \textbf{La jointure de fin de période}~: une seule ligne décide si le trimestre se terminant en $t$ apparaît en $t$ ou en $t+1$.
\omcode{code/ml/03-validation/python/ml_validation.py}{55}{63}{La surprise telle que la chaîne de traitement la voit, à la date de dépôt ou à la fin de période.}\label{lst:ml:validation:join}
\item \textbf{\omterm{def:ml:validation:nested}{Validation croisée imbriquée}} sur n'importe quel estimateur scikit-learn et n'importe quel découpeur.
\omcode{code/firm/cvsplit/firm_cvsplit.py}{70}{92}{Validation croisée imbriquée et meilleur score optimiste de la recherche à plat.}\label{lst:ml:validation:nested}
\item \textbf{Le test de troncature et le canari.}
\omcode{code/firm/cvsplit/firm_cvsplit.py}{102}{119}{Deux détecteurs de l'audit de fuite.}\label{lst:ml:validation:detectors}
\item \textbf{Lancer} \texttt{leaks()}, \texttt{overlap\_detectors()}, \texttt{detectors()}, \texttt{nested()} et
\texttt{fig\_validation.py}.
\end{tutsteps}
\textbf{Pour aller plus loin.} Remplacer l'encodage par la cible du même mois par un encodage calculé sur les rendements du mois précédent, et
vérifier que le test de troncature passe~; porter l'effet d'annonce à 4~\% et observer le R² fictif de la jointure de fin de période
croître avec son carré.
\end{tutorial}

\section{Construction~: découpeurs et détecteurs de fuites}\label{bld:ml:validation:cvsplit}

\begin{build}
\bfield{Rôle} Chaque modèle de l'entreprise est validé avec des découpages qui respectent le temps, choisis sans toucher à son
évaluation, et soumis à un \omterm{def:ml:validation:audit}{audit de fuite}.
\bfield{Interface} \texttt{PurgedKFold(t0, t1, n\_splits, embargo)},
\texttt{CPCVSplit(t0, t1, n\_folds, n\_test, embargo)}, \texttt{WalkForwardSplit(n, train, test, step, anchored)} (validateurs croisés scikit-learn, sur le
\texttt{firm.overfit} du livre~7)~; \texttt{nested\_cv(make\_model, grid, X, y, outer, inner, score)}~;
\texttt{fold\_overlap}, \texttt{truncation\_test}, \texttt{canary}, \texttt{adversarial\_validation}.
\bfield{Règles} Lignes en ordre chronologique~; chaque étiquette porte son intervalle~; un score utilisé pour choisir n'est jamais présenté comme une
évaluation.
\bfield{Tests d'acceptation} \texttt{code/firm/cvsplit/tests/}~: les découpeurs fonctionnent dans \texttt{cross\_val\_score} et
les plis purgés ne laissent aucun chevauchement~; la recherche imbriquée retient la bonne pénalité ridge, et le meilleur score à plat n'est jamais inférieur au
score imbriqué~; le test de troncature signale une moyenne mobile centrée et laisse passer la version causale~; le canari signale une chaîne qui
obtient un score sur une variable faite de sa cible~; la \omterm{def:ml:validation:audit}{validation adverse} est proche d'un demi sur une seule distribution et élevée sur deux.
\bfield{Pour aller plus loin} Un traceur de dépendances qui enregistre les lignes lues pour chaque valeur de variable~; un canari qui permute la cible
par blocs aussi longs que les étiquettes.
\end{build}

\begin{omsources}
\item S.~Varma et R.~Simon, ``Bias in error estimation when using cross-validation for model selection'', \emph{BMC
Bioinformatics} 7, 2006.
\item G.~C.~Cawley et N.~L.~C.~Talbot, ``On over-fitting in model selection and subsequent selection bias in performance
evaluation'', \emph{Journal of Machine Learning Research} 11, 2010.
\item S.~Kaufman, S.~Rosset, C.~Perlich et O.~Stitelman, ``Leakage in data mining: formulation, detection, and
avoidance'', \emph{ACM Transactions on Knowledge Discovery from Data} 6(4), 2012.
\item J.~Bergstra et Y.~Bengio, ``Random search for hyper-parameter optimization'', \emph{Journal of Machine Learning
Research} 13, 2012.
\end{omsources}

\section{Exercices}

\begin{exercise}[$\star$]\label{exo:ml:validation:1}
Cinq réglages d'\omterm{def:ml:why-financial-machine-learning-is-different:sets}{hyperparamètres} ont un \omterm{def:ml:why-financial-machine-learning-is-different:r2}{R² hors échantillon} véritable de 0.30~\% chacun, et chaque estimation de validation a
une erreur type de 0.10~\%, indépendamment. Quelle est à peu près la moyenne du meilleur score de validation~? (L'espérance du maximum de
cinq normales centrées réduites vaut 1.16.)
\end{exercise}

\begin{exercise}[$\star$]\label{exo:ml:validation:2}
Les étiquettes couvrent trois mois et sont échantillonnées chaque mois~; le pli de test va des mois 100 à 139. Avec une purge et un
embargo d'un mois, quels mois d'entraînement sont retirés~?
\end{exercise}

\begin{exercise}[$\star$]\label{exo:ml:validation:3}
Une exécution de canari d'une chaîne donne des scores de R² de 0.21~\%, 0.15~\% et 0.18~\% sur trois cibles de bruit. La chaîne
fuit-elle~? Calculer la moyenne et son erreur type.
\end{exercise}

\begin{exercise}[$\star\star$]\label{exo:ml:validation:4}
Pourquoi le canari manque-t-il la jointure en fin de période, alors que le test de troncature la détecte~? Lequel des deux détecte
l'encodage de la cible, et pourquoi faut-il les deux~?
\end{exercise}

\begin{exercise}[$\star\star$]\label{exo:ml:validation:5}
La \omterm{def:ml:validation:audit}{validation adverse} des mois d'entraînement et de test d'un modèle donne une AUC de 0.95. Donner deux conclusions différentes que cela
appuie, et une qu'il n'appuie pas.
\end{exercise}

\begin{exercise}[$\star\star$]\label{exo:ml:validation:6}
\emph{Trouver la faille.} «~Nous standardisons chaque variable avec la moyenne et l'écart type de l'ensemble du jeu de données, puis nous
lançons une validation croisée purgée, donc il n'y a aucune fuite.~»
\end{exercise}

\begin{exercise}[$\star\star\star$]\label{exo:ml:validation:7}
\emph{Programmation.} Dans \texttt{ml\_validation}, fixer l'effet d'annonce \texttt{ANN} à 4~\% et relancer les chaînes de fin de période et
de date de dépôt. De combien le R² factice augmente-t-il, et pourquoi environ quatre fois~?
\end{exercise}

\begin{exercise}[$\star\star\star$]\label{exo:ml:validation:8}
Démontrer que le score externe de la \omterm{def:ml:validation:nested}{validation croisée imbriquée} est un estimateur sans biais du score de généralisation de la procédure lorsque
les plis externes sont indépendants des choix de la procédure, et expliquer pourquoi le meilleur score à plat ne l'est pas.
\end{exercise}

\section{Problème~: Cinq fuites}

\begin{problem}[{Problème du week-end --- audit d'une chaîne qui ne peut pas gagner}]\label{pb:ml:validation:1}
Le monde du chapitre~: 200 actions, 240 mois, des caractéristiques qui ne portent aucune information, et des résultats trimestriels
publiés un mois après la fin de chaque trimestre.

\medskip\noindent\textbf{Partie I --- La chaîne honnête.}
\begin{enumerate}
\item Que calcule la chaîne soigneuse, et pourquoi un nombre négatif est-il ici la bonne réponse~?
\item Pourquoi retire-t-on un mois entre les mois d'entraînement et les mois de test~?
\item Que met la jointure sur la date de dépôt dans la ligne du mois $t$, et que met-y la jointure en fin de période~?
\item Quel effet d'annonce rend le R² factice de la jointure en fin de période égal à 1.27~\%~?
\end{enumerate}

\medskip\noindent\textbf{Partie II --- Les fuites.}
\begin{enumerate}[resume]
\item Donner le R² factice de l'encodage de la cible et expliquer sa taille.
\item Que donnent la sélection sur l'ensemble de l'échantillon et la sélection sur les mois d'entraînement~?
\item Combien d'étiquettes d'entraînement partagent un mois avec un pli de test mélangé, et combien avec un pli purgé~?
\item Que donnent les arbres profonds avec des plis mélangés et avec des plis purgés, et qu'est-ce qui leur permet d'exploiter le chevauchement~?
\end{enumerate}

\medskip\noindent\textbf{Partie III --- La sélection.}
\begin{enumerate}[resume]
\item Quelle plage couvrent les scores de sélection des trente graines, et à quelle distance de la médiane se trouve la meilleure~?
\item Que donne la graine retenue sur les mois non utilisés~?
\item Que rapportent les recherches à plat et imbriquée~?
\item Pourquoi la \cref{prop:ml:validation:max} prédit-elle l'écart~?
\end{enumerate}

\medskip\noindent\textbf{Partie IV --- L'audit.}
\begin{enumerate}[resume]
\item Quel détecteur attrape chaque fuite, avec son numéro~?
\item Pourquoi aucun détecteur seul ne peut-il attraper les cinq~?
\item Que signifie pour cette chaîne une AUC adverse de 0.71 sur les caractéristiques~?
\item Énoncer le \emph{résultat nommé}~: le R² que chaque fuite fabrique sur une cible qui ne peut pas être prédite, et le
détecteur qui la rattrape.
\item Rédiger l'audit sous forme de liste de contrôle qu'un relecteur signe.
\item Laquelle des cinq fuites survivrait à une période de trading simulé, et pendant combien de temps~?
\item Comment le journal de recherche doit-il consigner une fuite une fois trouvée~?
\item En une phrase~: à quoi sert la validation~?
\end{enumerate}
\end{problem}

\section{Questions d'entretien}

\begin{interviewq}[$\star$ \iqroles{chercheur, ingénieur ML}]\label{iq:ml:validation:1}
Qu'est-ce qu'une fuite de données~? Donnez trois exemples tirés de données financières.
\end{interviewq}

\begin{interviewq}[$\star\star$ \iqroles{ingénieur ML}]\label{iq:ml:validation:2}
Vous avez réglé vingt réglages d'\omterm{def:ml:why-financial-machine-learning-is-different:sets}{hyperparamètres} par validation croisée à 5 plis et rapportez le meilleur score moyen sur les plis. Qu'est-ce qui
cloche, et que rapportez-vous à la place~?
\end{interviewq}

\begin{interviewq}[$\star\star$ \iqroles{chercheur}]\label{iq:ml:validation:3}
Comment testeriez-vous si une chaîne de variables comporte un biais d'anticipation, sans lire son code~?
\end{interviewq}

\begin{interviewq}[$\star\star$ \iqroles{chercheur, ingénieur ML}]\label{iq:ml:validation:4}
Qu'est-ce que la \omterm{def:ml:validation:audit}{validation adverse}, et que vous apprend-elle sur une séparation entraînement--test~?
\end{interviewq}

\begin{interviewq}[$\star\star$ \iqroles{chercheur}]\label{iq:ml:validation:5}
Pourquoi l'encodage de la cible est-il dangereux avec des séries temporelles, et comment le faire en toute sécurité~?
\end{interviewq}

\begin{interviewq}[$\star\star\star$ \iqroles{chercheur}]\label{iq:ml:validation:6}
Le score croisé d'un modèle est excellent et son score en trading simulé est nul. Expliquez comment vous détermineriez si la cause est une fuite,
un \omterm{def:ml:why-financial-machine-learning-is-different:generalisation}{surapprentissage} ou un changement de marché.
\end{interviewq}
